\subsection{Extracción de características} %Falta incluir otros modelos que usemos y referencias
Para poder clasificar los audios en reales y generados tenemos que pensar primero en qué datos necesitamos pasar a los modelos para que aprendan. Podemos separar los datos que debemos proporcionar a dichos modelos en 3 grupos:

\begin{enumerate}
	\item Características. Podemos extraer características de la voz como el tono fundamental, los formantes, las pausas, y otros. Con estos datos, los modelos pueden aprender las relaciones entre los dos grupos de audios.
	\item Gráficas. Podemos extraer gráficas de la voz como la forma de onda, el espectrograma de MEL, y otros. Con estas gráficas las CNN pueden aprender los puntos clave que diferencian a los audios reales de los generados.
	\item Audio crudo. Existen modelos llamados wav2vec los cuales aprenden directamente de los audios, sacando un vector denso de características latentes en los audios, para después hacer ajuste fino o \emph{fine-tuning} de un modelo para clasificar los audios.  
\end{enumerate}

Vamos a ver más en detalle qué características de cada tipo sacamos.

\subsection*{Características}

Podemos extraer características de la voz como el tono fundamental, los formantes, las pausas, y otros. Con estos datos, los modelos pueden aprender las relaciones entre los dos grupos de audios. Cabe destacar que las características de este tipo extraídas son mayormente $arrays$ en el que cada componente es un valor en un momento determinado (series temporales). Esto se debe a que la voz tiene una componente temporal, por lo que hay que dividir cada audio en marcos temporales (frames). Además, los modelos que aplicaremos no entrenan con arrays, por lo que los promediaremos con la media y/o desviación típica. Vamos a explicar cada característica que extraemos de los audios:

\begin{enumerate}
	\item \textbf{Valor cuadrático medio (RMS, Root Mean Square)}: Mide la energía promedio de una señal acústica (i.e. la intensidad o volumen general del audio). Se extrae dividiendo la señal en frames, para cada frame se calcula la raíz cuadrada del promedio de los valores de onda al cuadrado y, por último, se promedian esos valores. En voz humana natural, el RMS tiende a tener variaciones más orgánicas, mientras que audios sintéticos pueden mostrar patrones más uniformes.
	\item \textbf{Tasa de cruces por cero (ZCR, Zero Cross Rate)}: Mide el número de veces que la señal cambia de signo por segundo, lo que nos permite distinguir entre los sonidos tonales (periódicos) y los ruidosos (no periódicos). Se extrae detectando los puntos donde la señal cambia de signo, se cuentan dichos puntos por unidad de tiempo y se normaliza por la longitud del frame. Los sonidos sordos (como /s/, /f/) tienen ZCR alto, mientras los sonoros (/a/, /m/) tienen ZCR bajo.
	\item \textbf{Centroide espectral}: Mide el centro de masa del espectro de frecuencias, lo que nos dice, de forma ponderada, si la energía de un sonido se concentra en las frecuencias graves (bajas) o en las frecuencias agudas (altas). Se obtiene calculando la Transformada Rápida de Fourier (FFT) para cada frame para obtener el espectro de frecuencias, se ponderan las frecuencias por su magnitud (energía de dicha frecuencia) y se calcula la media de estas. Cabe mencionar que las voces femeninas tienden a tener centroides más altos que las masculinas.
	\item \textbf{Ancho de banda espectral}: Mide qué tan dispersas están las frecuencias alrededor del centroide, para ver si el sonido es ``puro'' como un silbido o una sílaba mantenida, o es ``ruidoso'' como el ruido blanco o un aplauso. Se extrae hallando el centroide espectral y calculando la desviación típica ponderada en vez de la media ponderada. 
	\item \textbf{Punto de caída espectral (spectral rolloff)}: Mide la frecuencia por debajo de la cual se concentra un determinado porcentaje de la energía total del espectro, es decir, la frecuencia límite para un porcentaje de energía. Normalmente, dicho porcentaje es del 85\%, aunque puede variar entre 75\% a 90\%, y funciona como un percentil. Se obtiene calculando el centroide espectral y sumando la energía progresivamente hasta que lleguemos a una frecuencia con la que nos pasamos de dicho porcentaje. Es útil para distinguir entre sonidos con energía concentrada en bajas frecuencias, en un amplio espectro, o en altas frecuencias.
	\item \textbf{Planutid espectral}: Mide cómo de plano es el sonido. Varía entre 0 y 1 siendo 0 un sonido con picos pronunciados y 1 un sonido con pocos picos. Se obtiene calculando el espectro como antes (FFT) y dividiendo la media geométrica del espectro entre la media aritmética del mismo.
	\item \textbf{Tono fundamental (F0)}: Es la frecuencia más baja y periódica de un sonido. En el contexto de la voz, es la frecuencia a la que vibran las cuerdas vocales. El tono fundamental puede verse en, por ejemplo, las emociones, las preguntas frente a afirmaciones, las notas musicales, entre otras, por lo que nos da mucha información acerca de algo o alguien. Hay varias formas de calcularlo, y ninguna es 100\% precisa. Nosotros nos decantamos por la función PYIN, la cual primero obtiene las frecuencias fundamentales mediante el algoritmo YIN y después aplica el algoritmo probabilístico de Viterbi para obtener la cadena F0 más probable. Estos son algoritmos más sofisticados que combinan autocorrelación, análisis espectral y probabilidades para evitar errores comunes como la octava errónea (detectar 200Hz por ruido en vez del tono real de 100Hz), zonas no voceadas (unvoiced) y ruido de fondo.
	\item \textbf{Formantes}: Son las frecuencias de resonancia del tracto vocal (la boca, la garganta y la nariz). Estas frecuencias se amplifican de forma natural cuando hablamos, dependiendo de la forma que le demos a la boca.
	\begin{enumerate}
		\item F1 (Primer formante): Relacionado con la abertura de la mandíbula (qué tan abierta está la boca).
		\item F2 (Segundo formante): Relacionado con la posición de la lengua (adelante o atrás).
		\item F3 (Tercer formante): Relacionado con la posición de la punta de la lengua y tiene mucha importancia en las vocales y en la identidad del hablante.
	\end{enumerate}
Para extraerlos hay varias alternativas, nosotros usamos el método de Burg para estimar las formantes, después se muestrean 100 puntos a lo largo del audio y se calcula la media.
	\item \textbf{Relación Armónico-Ruido (HNR, Harmonic-to-Noise Ratio)}: Mide la proporción entre energía periódica (armónicos) y energía aperiódica (ruido). Esta relación es una indicadora directa de la eficiencia fonatoria y la salud vocal: puede indicar la existencia de patologías como voz ronca, o afonía; una fatiga vocal; o, en algunos casos, emociones. Hay varias formas de calcularla: como el método de autocorrelación o el de análisis espectral. Nosotros usamos la función to\_harmonicity de parselmouth.Sound, la cual la calcula con el método de autocorrelación: primero halla el F0 de los frames; después aplica la función de autocorrelación normalizada $r(\tau) = \frac{\int x(t) \cdot x(t+\tau) \, dt}{\int x^2(t) \, dt}$; luego busca el máximo de la función de autocorrelación en el intervalo alrededor del período fundamental esperado, este máximo corresponde a la correlación entre la señal y su versión desplazada un período; tras esto calcula el HNR en decibelios $\text{HNR} = 10 \cdot \log_{10}\left( \frac{r_{\text{max}}}{1 - r_{\text{max}}} \right)$ y, por último, aplica interpolación parabólica para obtener valores más precisos del máximo y realiza correcciones para compensar el efecto de la ventana de análisis. Es interesante saber que las voces humanas naturales suelen oscilar entre los 10 a 25 dB, por lo que unos valores extremos pueden indicar voz sintética.
	\item \textbf{Proporción de silencio}: Mide el porcentaje del tiempo o de frames donde no hay voz activa, es decir, hay silencio o ruido de fondo hasta cierto umbral. La forma de calcularla es muy simple: se establece un umbral de energía debajo del cual una energía es considerada silencio, se calcula el RMS y se le aplica el filtro del umbral. Finalmente se divide el número de frames considerados silencio entre el número total de frames. Esto, pese a su simpleza, puede ser determinante en el análisis de una voz ya que las IAs pueden generar patrones de pausas poco naturales o demasiado regulares.
	\item \textbf{Coeficientes Cepstrales en Frecuencia de Mel (MFCC)}: Son una representación compacta del espectro de frecuencias que imita cómo percibe el sonido el oído humano. En lugar de trabajar con frecuencias lineales (Hz), trabajan con una escala llamada escala Mel, que es más sensible a los cambios en frecuencias bajas que en altas (como nuestro oído). Son un total de 13 coeficientes: el primero representa la energía promedio del espectro (similar al RMS, pero en el dominio cepstral); mientras que el resto representan la forma detallada del espectro: los picos, los valles, las pendientes... Cada coeficiente añade un nivel de detalle diferente. Estos coeficientes son muy importantes ya que son, con diferencia, la característica más utilizada en: reconocimiento de voz (Speech-To-Text), reconocimiento de hablante (Identificación), clasificación de géneros musicales, entre muchas otras aplicaciones. Hoy en día son un estándar en reconocimiento de voz. Se extraen mediante un proceso más laborioso: primero se hace un pre-énfasis, amplificando las frecuencias altas, ya que estas tienen poca energía. Posteriormente, se calcula el espectro de frecuencias de cada frame; tras esto se aplica un banco de filtros triangulares espaciados según la escala Mel, donde por debajo de los 1000 Hz se aplican filtros lineales (nuestro oído distingue bien los cambios en frecuencias bajas) y por encima de los 1000 Hz se aplican filtros logarítmicos (nuestro oído distingue peor los cambios en frecuencias altas), con la fórmula $\text{Mel}(f) = 2595 \cdot \log_{10}\left(1 + \frac{f}{700}\right)$ y, para terminar, aplica la Transformada Coseno Discreta (DCT), la cual decorrelaciona las bandas de Mel, que suelen estar correlacionadas entre sí, compacta la información y devuelve los 13 primeros coeficientes (pueden aparecer entre 20 y 40, pero no son muy representativos).
	\item \textbf{Asimetría (skewness)}: Mide hacia qué lado se inclina la distribución de la energía en un conjunto de datos. Si hay una asimetría positiva, la energía se concentra en valores bajos, con una cola larga hacia los valores altos. En cambio si hay una asimetría negativa, la energía se concentra en valores altos, con una cola larga hacia los valores bajos. Si no hay asimetría, la energía está equilibrada alrededor del centro (como una campana de Gauss). Se calcula restando la media a cada muestra y elevando al cubo, luego hallando el promedio de estas desviaciones cúbicas, y, finalmente, dividiendo por la desviación estándar al cubo para estandarizar. Cabe mencionar que unas distribuciones de amplitud asimétricas pueden indicar características anómalas en voces sintéticas.
	\item \textbf{Curtosis}: Mide qué tan concentrados están los valores alrededor de la media y qué tan pesadas son las colas de la distribución. En audio, la curtosis nos dice qué tan ``espigado'' o ``suave'' es un sonido en términos de cómo se distribuyen sus amplitudes (en el tiempo) o su energía (en la frecuencia). Se obtiene con la siguiente fórmula $\text{Curtosis} = \frac{\frac{1}{N}\sum_{i=1}^N (x_i - \bar{x})^4}{\left(\sqrt{\frac{1}{N}\sum_{i=1}^N (x_i - \bar{x})^2}\right)^4}$. Adicionalmente, se puede restar 3 a la curtosis para que la mitad (mesocúrtica) se halle en 0. Podemos destacar que una curtosis alta (distribución picuda) puede aparecer en audios con compresión excesiva o artefactos de síntesis.
	\item \textbf{Transcripción}: Para la transcripción de los audios, utilizamos modelos de \href{https://alphacephei.com/vosk/models}{\textbf{Vosk}}. Esta elección se basó en que tiene modelos para los dos idiomas que vamos a investigar y además, dos versiones dentro de cada lenguaje. El que usamos para las pruebas es el small ya que es más rápido. Sobre la implementación, cabe mencionar que primero probamos Whisper, el modelo de reconocimiento automático del habla desarrollado por OpenAI, pero tras varios problemas durante su implementación, lo descartamos y escogimos Vosk en su lugar para esta tarea. Es interesante tener en cuenta la transcripción de un audio ya que en ocasiones, si no podemos sacar texto en claro, pero se sabe que hay voces, puede ser un primer indicio de que el audio es generado.
\end{enumerate}